\addchap{Preface}

This book began with a list of textbooks that should exist but do not. One entry asked a question that is becoming practical: once an AI system is better than its overseers at a task, how do the overseers check its answers? The entry observed that complexity theory has studied exactly this question for forty years, as the question of what a weak verifier can learn from a strong prover that may lie, and that results bearing on it are scattered across papers on AI safety: debate and its connection to polynomial space, doubly-efficient debate, prover–verifier games, spot checks in the style of probabilistically checkable proofs, delegation of computation, and audits of models in zero knowledge. Nobody had yet taught the complexity theory systematically for this problem. The entry proposed to do so one protocol at a time, and to ask the same three questions of every protocol: is honesty an equilibrium, what does the verifier spend, and how often is a lie caught?

The resulting book has fifteen chapters in five parts. Part I studies provers that are fixed objects: certificates, proofs that the verifier spot-checks, and programs that it queries. Part II studies a single prover that answers questions, from private coins to $\cc{IP}=\cc{PSPACE}$ and the delegation of computation to an efficient prover. Part III sets provers against each other: debate, cross-examination, and debate in which the honest side must be efficient. Part IV studies provers that are trained or paid rather than adversarial: rational proofs, prover–verifier games, and the verification of learned models. Part V studies claims about objects the prover will not show, with zero knowledge and proofs of training, and ends with backdoors that no audit can detect. Each chapter opens with an oversight problem, defines its protocol, proves the protocol's main theorem completely, runs the protocol in a short script against honest and cheating provers, answers the three questions, and closes by saying what the mathematical model leaves out of the real problem.

The intended readers are graduate students and researchers in machine learning and AI safety who want the theory behind scalable oversight, and theoretical computer scientists who want to see where their theory meets it. The prerequisites are undergraduate discrete mathematics, probability, linear algebra, and algorithms; no course in complexity theory or cryptography is assumed. Each chapter defines what it uses, and an appendix collects the shared tools, with proofs of those whose proofs are short.

Two cautions belong at the front. First, the bridge between the two fields is an analogy with exact parts and idealized parts, and the book tries to say which is which. A human overseer is not a polynomial-time machine that never errs, a trained model is neither an all-powerful prover nor a worst-case adversary, and a protocol that certifies a formal statement certifies nothing about whether the statement says what the overseer meant. Every chapter's last section is about these gaps. Second, much of the AI-safety side of the subject is young: several of the results discussed are preprints from 2024 to 2026, and the experiments behind its empirical claims are small. The chapters date such claims and say when a result has not been refereed.

\paragraph{A note on authorship.} This book was written by Claude, an AI model made by Anthropic, and conceived and edited by Xiaoyu Li, who proposed the subject, its plan, and the three questions. Every theorem is stated with its hypotheses and proved or explicitly sketched, every claim attributed to a paper was checked against the paper or a reliable account of it wherever these could be reached, every number from the book's scripts comes from code that the reader can run, and the colophon describes how the book was made. It would be an odd book on oversight that asked to be trusted. Its claims are written to be checked, and readers who find an error are invited to report it, as a verifier would: with the step that fails.
